\section{人格权}
\begin{enumerate}
  \item 1 道题目。\zd{识别权利主体、权利客体、侵权行为}。
  \item 停止侵害、排除妨碍、消除危险、消除影响、恢复名誉、赔礼道歉。
  \item 一般人格权。人身自由，\zd{非法拘禁}；人格尊严，\zd{非法搜查}。
  \item 具体人格权。
        \begin{enumerate}
          \item 物质性人格权。
                \begin{enumerate}
                  \item 生命。得死。
                  \item 身体。\zd{完整性和行动自由}。\zd{能否自由拆除}，\al{假牙，假肢}。
                  \item 健康。生理、心理\zd{机能}（功能）的正常运转和发挥。
                \end{enumerate}
          \item 精神性\blh
                \begin{enumerate}
                  \item 姓名/名称。扩大解释，笔名、艺名。
                        \begin{enumerate}
                          \item 干涉。不允许别人用。
                          \item \zd{盗用}。只用姓名不用身份。
                          \item 假冒。身份 + 姓名。
                        \end{enumerate}
                  \item 肖像。可以\zd{被识别的外部形象}，扩大解释。
                        \begin{enumerate}
                          \item 丑化，\al{葛优躺}；污损；伪造，\al{PS，AI 换脸}；制作，\al{表情包}；使用，\al{用别人照片当头像}；公开。
                          \item 合理使用制度。实时新闻报道；职务行为（\al{逃犯画像}）；寻人启事。
                        \end{enumerate}
                  \item \sjt 许可使用。\al{明星代言}。
                  \item 名誉。\zd{贬低；负面社会评价}（含\zd{信用}）。侮辱；诽谤（捏造散布虚\zd{假}事实）。文学作品真人真事为对象算侵权。
                  \item 荣誉。光荣称号。剥夺、诋毁、贬损。\al{邱少云案；袁隆平去世}。
                  \item 隐私。常考。私密空间、活动、信息、\zd{私人生活安宁}；\zd{刺探、侵扰、泄露、公开}。
                        \begin{enumerate}
                          \item \al{骚扰电话；垃圾短信；针孔摄像头；人肉搜索；私家侦探；狗仔}。
                        \end{enumerate}
                  \item 婚姻自主。从来不考。
                \end{enumerate}
        \end{enumerate}
  \item 个人信息保护。\zd{可识别到特定人}。\al{脸、身份证、指纹、基因}；\zd{非法收集、存储、传输、买卖}。
        \begin{enumerate}
          \item 和数据区别。无可识别性。
          \item 和隐私。个人信息是大圈，隐私是小圈；姓名是个人信息，但不是隐私。
          \item 免责。为了社会公共利益。
        \end{enumerate}
\end{enumerate}
